\section{QED und die Rolle der Masselosigkeit für die Renormierbarkeit}

Die Quantenelektrodynamik (QED) gehört zu den präzisesten und erfolgreichsten 
Theorien der modernen Physik. Ein zentraler Grund für ihre interne 
Konsistenz ist die Masselosigkeit des Photons. Diese Eigenschaft ist nicht 
nur eine empirische Annahme, sondern eine strukturelle Notwendigkeit der 
Theorie. Eine – selbst extrem kleine – Photonenmasse würde die 
Eichsymmetrie brechen und die Renormierbarkeit sowie die mathematische 
Kohärenz der QED zerstören.

\subsection*{Eichsymmetrie als Fundament der QED}

Die QED beruht auf der lokalen $U(1)$-Eichsymmetrie des elektromagnetischen 
Feldes. Diese Symmetrie legt fest:

\begin{itemize}
	\item die Form der Kopplung zwischen Photonen und geladenen Teilchen,
	\item die Struktur des elektromagnetischen Stroms,
	\item die Erhaltung der elektrischen Ladung,
	\item und die notwendigen Beziehungen, durch die Ultraviolettdivergenzen 
	wegfallen.
\end{itemize}

Die Ward–Takahashi-Identitäten, direkte Konsequenzen der Eichsymmetrie, 
stellen sicher, dass Propagatoren, Vertexfunktionen und Schleifenkorrekturen 
die konsistenten Relationen behalten, die für eine renormierbare Theorie 
erforderlich sind.

Ein Photonenmasseterm bricht die Eichsymmetrie und macht diese Identitäten 
ungültig.

\subsection*{Zerfall der Ward–Takahashi-Identitäten}

In der masselosen QED erzwingen die Ward–Takahashi-Identitäten Relationen 
der Form
\[
q_\mu \Gamma^\mu = S^{-1}(p+q) - S^{-1}(p),
\]
wobei $\Gamma^\mu$ die Vertexfunktion und $S(p)$ der Fermionpropagator ist.  
Diese Beziehung garantiert die systematische Aufhebung von Divergenzen und 
schützt die Struktur der elektrischen Ladung.

Ein Photonenmasseterm verletzt genau die Symmetrie, aus der diese Identitäten 
folgen. Dadurch:

\begin{itemize}
	\item heben sich Divergenzen nicht mehr gegenseitig auf,
	\item wird die Ladungsrenormierung inkonsistent,
	\item verliert die Theorie ihre Vorhersagekraft.
\end{itemize}

Damit ist klar: Eine renormierbare QED erfordert zwingend $m_\gamma = 0$.

\subsection*{Longitudinale Moden und nicht-renormierbare Beiträge}

Ein massives Photon besitzt zusätzlich zur transversalen Polarisation einen 
longitudinalen Freiheitsgrad. In Streuprozessen, an denen dieser Modus 
beteiligt ist, wachsen die Amplituden in der Regel mit der Energie.  
In einer eichinvarianten Theorie werden solche Wachstumsbeiträge durch 
Interferenzen zwischen Diagrammen kompensiert. Ohne Eichsymmetrie fallen 
diese Auslöschungen weg.

Die Folgen sind:

\begin{itemize}
	\item Verlust der perturbativen Unitarität,
	\item Auftreten nicht-renormierbarer Divergenzen,
	\item Verletzung fundamentaler Wahrscheinlichkeitsgrenzen.
\end{itemize}

Damit wird deutlich: Masselosigkeit ist keine Option, sondern eine 
Voraussetzung für die quantenkonsistente Beschreibung des 
elektromagnetischen Feldes.

\subsection*{Verbindung zur elektroschwachen Theorie}

Im Standardmodell entsteht das Photon nach dem spontanen Symmetriebruch aus 
der elektroschwachen Eichstruktur. Seine Masselosigkeit ist durch die 
ungebrochene elektromagnetische $U(1)$-Symmetrie geschützt.  
Im Gegensatz dazu erhalten die $W^\pm$- und $Z$-Bosonen ihre Massen durch 
den Higgs-Mechanismus.

Ein Photonenmasseterm würde einen Mechanismus erfordern, der über das 
Standardmodell hinausgeht – einen Mechanismus, der notwendigerweise die 
elektromagnetische Eichsymmetrie bricht und damit die renormierbare Struktur 
der QED zerstört.

\subsection*{Konzeptionelle Bedeutung}

Die Masselosigkeit des Photons ist tief in die Quantenstruktur der 
Elektrodynamik eingebettet:

\begin{itemize}
	\item sie garantiert die Renormierbarkeit,
	\item sie stellt die Erhaltung der elektrischen Ladung sicher,
	\item sie stabilisiert das Hochenergieverhalten von Streuamplituden,
	\item sie ermöglicht die störungstheoretischen Auslöschungen, die erst 
	durch Eichsymmetrie möglich werden.
\end{itemize}

Diese Eigenschaften sind nicht verhandelbar, sondern essenziell für die 
mathematische und physikalische Konsistenz der Theorie. Ein massives Photon 
würde daher eine vollständig neue, alternative Feldtheorie der 
Elektromagnetismus erfordern.

\subsection*{Übergang zur Rolle von Helizität und Verschränkung}

Das nächste Kapitel wendet sich einem experimentell besonders klaren Fingerabdruck 
der Photonenmasselosigkeit zu: der Helizitätsstruktur verschränkter Photonen.  
Dass Photonen nur zwei Helizitätszustände besitzen, ist sowohl theoretisch als 
auch experimentell ein eindeutiges Merkmal eines masselosen Spin-1-Teilchens.  
Ein massives Photon hätte einen dritten, longitudinalen Zustand, der die 
Struktur von Verschränkungs- und Bell-Experimenten messbar verändern würde.

Damit verbinden sich die theoretischen Argumente der QED direkt mit 
beobachtbaren quantenmechanischen Phänomenen.
